
\begin{abstract}
% 5-7% of total word count (~150-210 words for a 3,000-word paper)
[1-2 sentences: Empirical puzzle/hook]
[1-2 sentences: Motivation/importance]
[1 sentence: Research question]
[1-2 sentences: Research design/approach]
[1-2 sentences: Main findings]
[1 sentence: Contribution/implications]

\medskip
\noindent\textbf{Keywords:} [Keyword 1], [Keyword 2], [Keyword 3], [Keyword 4], [Keyword 5]
\end{abstract}

\newpage

% =====================================================================
\section{Introduction}\label{sec:intro} % 8-10% of total (~240-300 words)

\textbf{[Opening puzzle / hook]} % ~20% of intro
[A striking empirical puzzle, anomaly, or paradox. One powerful observation. Avoid generic openings.]

\textbf{[Motivation \& `so what?']} % ~20% of intro
[(1) Why does this matter for the real world? (2) Why does it matter for political science theory? Address head-on in 2-3 sentences.]

\textbf{[Research question]} % ~10% of intro
[One single, clear, falsifiable question. Prefer: `Under what conditions / To what extent / Why does X affect Y in context Z?']

\textbf{[Literature gap - brief]} % ~15% of intro
[In 3-4 sentences: what do we know (3-5 key works), and what is missing? End with: `This article addresses this gap by...']

\textbf{[Main argument / thesis]} % ~10% of intro
[One to two sentences stating your answer to the RQ. Be direct and specific. No hedging verbs like `suggest' or `explore'.]

\textbf{[Contribution statement]} % ~15% of intro
[What this article adds: new case, mechanism, test, dataset, or methodological advance. Use: `I contribute to the literature on X by Y.']

\textbf{[Roadmap]} % ~10% of intro
[One sentence per section, cross-referenced with \texttt{\textbackslash cref}:\Cref{sec:lit} reviews\dots{} \Cref{sec:theory} develops\dots{} \Cref{sec:design} presents\dots{} \Cref{sec:results} reports\dots{} \Cref{sec:discussion} discusses\dots]

% =====================================================================
\section{Literature Review}\label{sec:lit} % 15-20% of total (~450-600 words)

\subsection{Theoretical Foundations} % 25% of lit review
[Place your article within a recognised theoretical tradition (rationalism, constructivism, institutionalism, etc.). Show intellectual lineage and explain why this tradition is the relevant starting point.]

\subsection{First Debate: [Theme 1]} % 30% of lit review
[Identify the first major strand. Synthesize, do NOT list papers one by one. State: central claim $\rightarrow$ key works (3-5 citations, e.g.\ \citet{author2020} or \citep{author2020}) $\rightarrow$ empirical scope $\rightarrow$ limitations.]

\subsection{Second Debate: [Theme 2]} % 30% of lit review
[Identify the second strand (or contrasting school). Same structure: central claim $\rightarrow$ key works $\rightarrow$ limitations. Make clear how it relates to or competes with Debate 1.]

\subsection{The Gap and Transition to Theory} % 15% of lit review
[Synthesize the state of the field into a single, precise gap statement. Use: (a) untested mechanism, (b) neglected context, (c) contradictory findings, (d) missing variable.]
[Explicit bridge: `The gap identified above calls for a theoretical account of [mechanism]. The following section provides such an account.']

% =====================================================================
\section{Theoretical Framework}\label{sec:theory} % 15-20% of total (~450-600 words)

\subsection{Core Concepts and Definitions} % 15% of theory
[Define your key theoretical concepts precisely. Operational definitions come later in Research Design.]

\subsection{Causal Mechanism} % 35% of theory
[Describe step-by-step how X causes Y, using an explicit diagram (\cref{fig:dag}). Be explicit about each link. Avoid vague language like `X influences Y'; say how and why each step operates.]

\begin{figure}[htbp]
    \centering
    \footnotesize
    \begin{tikzpicture}[box/.style={draw, rounded corners, align=center, minimum width=3.5cm, minimum height=1.2cm}]
        \node[box] (X)  at (0,0)  {D: [Treatment]};
        \node[box] (M1) at (5,3)  {$M_1$: [Mechanism 1]};
        \node[box] (M2) at (5,1)  {$M_2$: [Mechanism 2]};
        \node[box] (M3) at (5,-1) {$M_3$: [Mechanism 3]};
        \node[box] (M4) at (5,-3) {$M_4$: [Mechanism 4]};
        \node[box] (U)  at (5,5.5) {U: [Pre-treatment\\characteristics]};
        \node[box] (Y)  at (10,0) {Y: [Outcome]};

        \foreach \m in {M1, M2, M3, M4} \draw[-{Stealth}, thick] (X) -- node[above, sloped] {$+$} (\m);
        \foreach \m in {M1, M2, M3}     \draw[-{Stealth}, thick] (\m) -- node[above, sloped] {$+$} (Y);
        \draw[-{Stealth}, thick] (M4) -- node[above, sloped] {$-$} (Y);

        \draw[-{Stealth}, dashed] (U) -- (X);
        \draw[-{Stealth}, dashed] (U) -- (Y);
    \end{tikzpicture}
    \caption{[DAG: causal pathways from X to Y]}
    \label{fig:dag}
\end{figure}

\subsection{Theoretical Predictions} % 35% of theory

\textbf{Main Hypothesis (H1)}
[One falsifiable, directional statement: `As X increases, Y increases/decreases, holding Z constant.' Write in plain language AND as a formal if-then statement.]

\textbf{Alternative Hypothesis / Rival Explanation (H2)}
[What would competing theories predict? State at least one rival hypothesis. Explain why your theory predicts a different outcome.]

\textbf{Conditional Hypothesis (H3, if relevant)}
[If your argument is conditioned on a third variable Z: `The effect of X on Y is stronger when Z is high.']

\subsection{Scope Conditions} % 15% of theory
[Define the population of cases for which your theory holds. What are the necessary background conditions? Explicit scope conditions protect against over-generalization.]

\subsection{Empirical Implications} % 15% of theory
[State explicitly what observable implications follow from your theory. What should we observe in the data if H1 is true vs.\ if H2 is true? This bridges theory and research design.]

% =====================================================================
\section{Research Design and Data}\label{sec:design} % 15-20% of total (~450-600 words)

\subsection{Causal Identification Strategy} % 20% of design
[Explain the core research design: OLS/logit, DiD, RDD, IV, matching, etc. Justify the choice. State the identifying assumption clearly and discuss why it is plausible in your context.]

\subsection{Sample, Population, and Time Period} % 12% of design
[Unit of analysis? N? Years? Justification for temporal/geographic scope? Selection concerns?]

\subsection{Dependent Variable (Y)} % 12% of design
[Operationalization? Source? Unit of measurement? Validity (does it capture the theoretical concept)? Name the variable exactly as in the analysis.]

\subsection{Independent Variables (X)} % 12% of design
[Same structure as Y. Is there endogeneity risk? How do you address it?]

\subsection{Control Variables} % 12% of design
[List each control with a one-sentence theoretical justification. Avoid kitchen-sink controls: only include variables that theoretically confound the X--Y relationship.]

\subsection{Descriptive Statistics} % 12% of design
[Reference the table explicitly (\cref{tab:descr}). Note any unusual distributions or missingness.]

% tip: once the R scripts write the table, replace the tabular below with
% \input{../4_outputs/tables/descriptives.tex}
\begin{table}[htbp]
    \centering
    \footnotesize
    \caption{Descriptive statistics}
    \label{tab:descr}
    \begin{tabular}{l*{5}{c}}
        \toprule
        Variable & N & Mean & SD & Min & Max \\
        \midrule
        {[Variable 1]} & [N] & [Mean] & [SD] & [Min] & [Max] \\
        {[Variable 2]} & [N] & [Mean] & [SD] & [Min] & [Max] \\
        {[Variable 3]} & [N] & [Mean] & [SD] & [Min] & [Max] \\
        \bottomrule
    \end{tabular}
\end{table}

\subsection{Model Specification} % 20% of design
[Write out the estimating equation (\cref{eq:main}). Explain the choice of standard errors (clustered, robust, bootstrapped). If panel data: fixed or random effects and why.]

\begin{equation}
    Y_{it} = \alpha + \beta X_{it} + \gamma' Z_{it} + \mu_i + \tau_t + \varepsilon_{it}
    \label{eq:main}
\end{equation}

% =====================================================================
\section{Results}\label{sec:results} % 15-20% of total (~450-600 words)

\subsection{Main Results} % 40% of results
[Present the main regression table (\cref{tab:results}). Discuss substantive magnitude, not just statistical significance. Translate the coefficient: `A one-unit increase in X is associated with [effect size], which represents [\% of mean / SD units / real-world benchmark].']

\begin{table}[htbp]
    \centering
    \begin{threeparttable}
        \caption{Main results}
        \label{tab:results}
        \begin{tabular}{l*{3}{c}}
            \toprule
             & (1) & (2) & (3) \\
            \midrule
            {[IV]}       & [Coef] & [Coef] & [Coef] \\
                         & ([SE]) & ([SE]) & ([SE]) \\
            {[Controls]} & [Coef] & [Coef] & [Coef] \\
                         & ([SE]) & ([SE]) & ([SE]) \\
            \midrule
            N            & [N]    & [N]    & [N]    \\
            $R^2$        & [$R^2$] & [$R^2$] & [$R^2$] \\
            \bottomrule
        \end{tabular}
        \begin{tablenotes}[flushleft]
            \footnotesize
            \item \textit{Note:} [Explain model specifications. Standard errors in parentheses, clustered by [unit]. $^{*}p<0.1$; $^{**}p<0.05$; $^{***}p<0.01$.]
        \end{tablenotes}
    \end{threeparttable}
\end{table}

[Include at least one figure (coefficient plot, marginal effects, or discontinuity plot) to visualize the main result (\cref{fig:main}).]

\begin{figure}[htbp]                 % floating box; [htbp] = here, top, bottom, own page
    \centering
    \safegraphic[width=0.8\textwidth]{main_effect.pdf}  % 80% of the text width
    \caption{[Visualization of main effects]}            % figure captions go below
    \label{fig:main}                                     % cite it with \cref{fig:main}
\end{figure}

\subsection{Tests of Alternative Hypotheses} % 25% of results
[Show that H2 (rival explanation) receives less empirical support than H1: separate regression including the rival variable, placebo test, mechanism test, or subsample analysis.]

\subsection{Robustness Checks} % 20% of results
[Test the main result under: (1) alternative operationalization of X or Y, (2) alternative model specification, (3) restricted sample, (4) alternative bandwidth/caliper. State briefly that results hold, or explain if they do not.]

\subsection{Placebo and Falsification Tests} % 15% of results
[Summary of placebo tests (e.g., fake cutoff points, pre-treatment outcomes), see \cref{fig:grid}. These strengthen causal claims and address identification concerns. Full results in the appendix.]

\begin{figure}[htbp]
    \centering
    \begin{subfigure}{0.48\textwidth}   % 48% wide, so 2 fit per row
        \safegraphic{placebo_a.pdf}
        \caption{[First]}
    \end{subfigure}\hfill               % \hfill pushes the 2nd one to the right
    \begin{subfigure}{0.48\textwidth}
        \safegraphic{placebo_b.pdf}
        \caption{[Second]}
    \end{subfigure}

    \medskip                            % vertical space between rows

    \begin{subfigure}{0.48\textwidth}
        \safegraphic{placebo_c.pdf}
        \caption{[Third]}
    \end{subfigure}\hfill
    \begin{subfigure}{0.48\textwidth}
        \safegraphic{placebo_d.pdf}
        \caption{[Fourth]}
    \end{subfigure}
    \caption{[Placebo tests in a 2$\times$2 grid]}
    \label{fig:grid}
\end{figure}

% =====================================================================
\section{Discussion and Conclusion}\label{sec:discussion} % 10-12% of total (~300-360 words)

\subsection{Summary of Argument and Findings} % 20% of discussion
[Briefly restate: (1) research question, (2) theoretical argument, (3) main empirical finding. 3-4 sentences. Introduce no new evidence.]

\subsection{Interpretation of Findings} % 25% of discussion
[What do the results tell us about the theoretical mechanism? Do they support or challenge existing theories? What do null results (if any) mean? Be precise.]

\subsection{Theoretical Implications} % 20% of discussion
[What must we revise, nuance, or add to the existing literature? Name specific works your results corroborate or challenge. Be precise about what your contribution changes.]

\subsection{Limitations} % 15% of discussion
[Be honest: external validity, measurement limitations, remaining endogeneity threats, scope condition boundaries. Naming limitations is scholarly maturity, not weakness.]

\subsection{Future Research Agenda} % 15% of discussion
[What are the logical next steps? Name 2-3 specific open questions that follow from your findings. This situates your work as part of a research program.]

\subsection{Closing Statement} % 5% of discussion
[Return to the opening puzzle from the introduction and resolve it in light of your findings. One to two memorable sentences. Avoid cliches like `More research is needed.']

\newpage
\printbibliography

\newpage
\appendix
\section{Variable Definitions and Coding Rules}\label{app:variables}
[Detailed definitions, sources, and coding rules for all variables.]

\section{Institutional Background}\label{app:background}
[Contextual information about the case, if applicable.]

\section{Additional Figures and Tables}\label{app:additional}
[Supplementary analyses, balance tests, bandwidth sensitivity, etc.]

\section{Data Availability and Replication Code}\label{app:replication}
[Repository: {{repo_url}}. Run \texttt{renv::restore()}, then the scripts in \texttt{2\_scripts/} in numeric order, then compile this file.]

\end{document}
